% Part: sets-functions-relations
% Chapter: sets
% Section: basics

\documentclass[../../../include/open-logic-section]{subfiles}

\begin{document}

\olfileid{sfr}{set}{bas}
\olsection{Extensionaliteit: dezelfde elementen, dezelfde verzameling}

Een \emph{verzameling} bestaat uit dingen die we samen als één ding
bekijken. De dingen die erin zitten, heten \emph{elementen} of
\emph{leden} van de verzameling. Als $x$ !!a{element} van een
verzameling~$A$ is, schrijven we $x \in A$. Zo niet, dan schrijven we
$x \notin A$. Een verzameling zonder !!{element}s heet de
\emph{lege} verzameling. We schrijven die als~‘$\emptyset$’.

\begin{explain}
Het maakt niet uit hoe we de verzameling \emph{omschrijven}, in welke
\emph{volgorde} we haar !!{element}s opschrijven of hoe \emph{vaak}
we dezelfde !!{element}s meetellen. Alleen welke !!{element}s erin
zitten, telt. Dat
vatten we samen in de volgende regel.
\end{explain}

\begin{defn}[Extensionaliteit]
  Als $A$ en $B$ verzamelingen zijn, dan geldt $A = B$ precies als elk
  !!{element} van~$A$ ook !!a{element} van~$B$ is, en omgekeerd.
\end{defn}

Door extensionaliteit mogen we een handige schrijfwijze gebruiken.
Neem dingen $a_{1}$, \dots, $a_{n}$. Dan is
$\{a_{1}, \dots, a_{n}\}$ \emph{de} verzameling met de !!{element}s
$a_1, \ldots, a_n$. Het woord ‘\emph{de}’ is belangrijk: de regel van
extensionaliteit zegt dat er maar \emph{één} zo'n verzameling kan
zijn. Daarom geldt ook:
  \[
    \{a, a, b\} = \{a, b\} = \{b,a\}.
  \] 
Bij een verzameling maakt de volgorde van de !!{element}s dus niet
uit. Het maakt ook niet uit hoe vaak we hetzelfde element opschrijven.

\begin{tagblock}{novice}
\begin{ex}
Je kunt dingen bij elkaar nemen en daar een verzameling van maken.
Neem bijvoorbeeld de broers en zussen van Richard. Dat is één
persoon, dus we kunnen die verzameling schrijven als
$S=\{\textrm{Ruth}\}$. De positieve gehele getallen onder $4$ vormen
de verzameling $\{1, 2, 3\}$. Je mag die ook schrijven als
$\{3, 2, 1\}$, of zelfs als $\{1, 2, 1, 2, 3\}$. Volgens de regel van
extensionaliteit is dat steeds dezelfde verzameling. Elk !!{element}
van $\{1, 2, 3\}$ is namelijk ook !!a{element} van $\{3, 2, 1\}$ (en
van $\{1, 2, 1, 2, 3\}$), en omgekeerd.
\end{ex} 
\end{tagblock}

Vaak zeggen we welke dingen in een verzameling zitten door een
eigenschap te noemen die de !!{element}s allemaal hebben. Daar is een korte
schrijfwijze voor: $\Setabs{x}{\phi(x)}$. Hier staat $\phi(x)$ voor
de eigenschap die~$x$ moet hebben om bij de !!{element}s van de
verzameling te horen.

\begin{tagblock}{novice}
\begin{ex}
In ons voorbeeld hadden we $S$ ook zo kunnen schrijven:
\[
S = \Setabs{x}{x \text{ een broer of zus van Richard is}}.
\]
\end{ex}
\end{tagblock}

\begin{tagblock}{math}
\begin{ex}
Een getal heet \emph{volmaakt} precies als het gelijk is aan de som
van zijn echte delers. Dat zijn de getallen waar je het zonder rest door
kunt delen, behalve het getal zelf. Zo is $6$ volmaakt: zijn echte
delers zijn $1$, $2$ en~$3$, en $6 = 1 + 2 + 3$. Het getal $6$ is
zelfs het enige positieve gehele getal onder $10$ dat volmaakt is.
Volgens de regel van extensionaliteit mogen we daarom schrijven:
\[
	\{6\} = \Setabs{x}{x\text{ volmaakt is en }0 \leq x \leq 10}
\]
De rechterkant lezen we als ‘de verzameling van alle $x$ waarvoor $x$
volmaakt is en $0 \leq x \leq 10$’. Deze gelijkheid laat nog eens zien
dat het niet uitmaakt hoe we een verzameling omschrijven. De regel van
extensionaliteit zorgt er in het algemeen voor dat er nooit meer dan één
verzameling van alle $x$ kan zijn waarvoor $\phi(x)$ geldt. Daarom
mogen we $\Setabs{x}{\phi(x)}$ \emph{de} verzameling van alle $x$
noemen waarvoor~$\phi(x)$ geldt.
\end{ex}
\end{tagblock}

De regel van extensionaliteit geeft ons een manier om te bewijzen dat
twee verzamelingen gelijk zijn. Om $A = B$ te bewijzen, laat je zien:
als $x \in A$, dan $x \in B$; en als $y \in B$, dan $y \in A$.

\begin{prob}
Bewijs dat er hoogstens één lege verzameling is. Laat dus zien dat als
$A$ en $B$ verzamelingen zonder !!{element}s zijn, dan $A = B$.
\end{prob}
\end{document}
